\clearpage
\chapter{KODE PROGRAM}
\label{chap:kode_program}

Implementasi kode program dalam penelitian ini disusun menggunakan bahasa pemrograman Python pada platform Google Colab. Cakupan kode ini meliputi tahapan prapemrosesan data, rekayasa fitur (\textit{feature engineering}), pendekatan model hibrida \textit{Random Forest}-ARIMA, penyetelan parameter menggunakan Optuna, serta analisis kontribusi fitur menggunakan SHAP.

Dokumentasi mengenai langkah implementasi, konfigurasi model, serta repositori hasil komputasi dapat diakses secara terbuka melalui tautan berikut:

\begin{center}
    \small\url{https://colab.research.google.com/drive/12TNBPOXTBSwVDy9xRgqy1A4GX0xdm2Z1?usp=sharing}
\end{center}

Beberapa bagian kode program yang berkaitan langsung dengan penjelasan teknis juga disertakan pada Bab IV dan Bab V sebagai pelengkap untuk membantu memperjelas alur analisis penelitian.